%!TEX root = groundunitbriefing.tex

% LTeX: enabled=true

% LTex: dictionary+=SAMs
% LTex: dictionary+=CCUs
% LTex: dictionary+=Oerlikon
% LTex: dictionary+=Skyguard
% LTex: dictionary+=APCs
% LTex: dictionary+=IFVs
% LTex: dictionary+=Starstreak
% LTex: dictionary+=Mitla
% LTex: dictionary+=Landrovers
% LTex: dictionary+=AMLs
% LTex: dictionary+=Bv
% LTex: dictionary+=Panhard
% LTex: dictionary+=FACs

\chapter{Additional Rules}

\rulesection{Ground Units}

A ground \emph{unit} is a combat or logistics formation of infantry, armor, transports, AAA, SAM launchers, or radars. This section discusses them in general.%
\note{Rules}{
    Much of this section should be included in the 3A rules, but it's here for the moment.
}

\rulesubsection{Type and Size}

A ground unit will have a type, for example, “heavy tank” or “ZSU-23-4”, and a size, which will be one of squad, section, platoon, or battery. Typically, a squad has five to ten soldiers and/or one vehicle, gun, or launcher; a section has ten to twenty soldiers and/or two vehicles, guns, or launchers; and a platoon or battery has thirty to fifty soldiers and/or three to six vehicles, guns, or launchers.
The term “platoon” is more commonly used for infantry and tank units and “battery” for artillery and air-defense units, but within the game they are equivalent.

Companies, battalions, and other larger units above platoons or batteries are represented by multiple units. For example, an infantry company might be represented by three infantry platoon units.

All ground-combat and transport platoons are also available as sections and squads and all AAA platoons and batteries are also available as sections.
A section or squad is identical in every way to a platoon or battery that has taken 1D or 2D damage, respectively, except that no VPs are awarded for this notional damage and non-AAA squads and sections may not engage in barrage fire.%
\note{Squads and Sections as Damaged Platoons}{
    So for example sections/squads have 2D/1D damage resilience and \plus{1}/\plus{2} modifiers to AAA attacks. Similar modifiers should be applied to ground-to-ground attacks.
}

In the game, ground units are uniformly referred to as squads, sections, platoons, and batteries. However, other terms are used for these units even within English-speaking armies.
For example, in many Commonwealth armies a section is equivalent to a squad, a troop is equivalent to a platoon, and a squadron is equivalent to a company.
In the US Armored Cavalry, a troop is equivalent to a company and a squadron is equivalent to a battalion.

\rulesubsection{Basic Properties}

The type and size of a ground unit determines its basic properties: damage resilience, defense strength, target class, mobility, and sighting range.

%Many units also have properties related to their air-defense capabilities.

%\begin{itemize}
%    \item basic type (e.g., infantry or tanks);
%    \item size;
%    \item damage resilience;
%    \item defense strength and target class;
%    \item mobility;
%    \item sighting range in hexes;
%    \item the VPs awarded for 3D/2D/D damage.
%\end{itemize}

The damage resilience is the amount of damage a ground unit can suffer before being killed. It is determined by the size of the unit: a squad has a damage resilience of 1D; a section has a damage resilience of 2D; and a platoon or battery has a damage resilience of 3D.%
\note{Damage Resilience in {\AirStr}}{
    There is a form of damage resilience in {\AirStr}, in that units whose counters have white text are only able to take 1D of damage before being killed (see the last paragraph of rule 28). Such units include infantry SAMs, armored CCUs, and many vehicle SAMs. {\TSOH} also mentions units being able to take only 2D damage, but it doesn't have any relevant units since it attaches SAM teams to infantry platoons rather than having them as separate (see below) and doesn't have a separate FAC unit. In both cases, the idea seems to have been incompletely implemented in the rules. I'm generalizing and homogenizing the idea. This is important as we need sections for mobile AAA units with two vehicles and for the Oerlikon GDF (as Skyguard can only control two guns and not a full battery). To a large degree, a section is equivalent to a platoon that has taken 1D and a squad is equivalent to a platoon that has taken 2D, so the mechanics of the game are really not impacted.
}

The numerical defense strength is used to calculate the odds for air-to-ground attacks and for ground-to-ground combat.

The target class is used to determine the appropriate attack strength of any weapons used in air-to-ground attacks.
The class can be a hard, open, or soft according to the degree to which the unit enjoys armored protection.
Many weapons have different attack strengths against targets of different classes.
A hard target has armored protection from both direct (front, side, and rear) and top attacks.
Hard targets are further described as light, medium, or heavy according to their degree of protection, and have corresponding defense strengths.%
\note{Hard Targets}{
    Light, medium, and heavy armor correspond to defense strengths of \underline{4}, \underline{6}, and \underline{8} in {\AirStr}.
}
An open target has armored protection from direct attacks but not from top attacks.
An open target is normally treated as a hard target, but against napalm, fire, FAE, and air-burst HE bombs, it is considered to be a soft target.%
\note{Open Targets}{
    The BTR-50P APC and the M3 half-tracked APC are examples of open targets. In the 1956 and 1967 Arab-Israeli Wars where both sides used M3 half-tracks. If an M3 is attacked by an air-burst or napalm bomb, the crew and passengers will suffer more or less as if they were in the open.
}
A soft target does not have armored protection their personnel and/or weapons.%
\note{Soft Targets}{
    Examples of soft targets include infantry, towed artillery, unarmored vehicles, and armored vehicles that do not fully protect their personnel and/or weapons such as the M-107 and M-110 tracked artillery vehicles.

    Nothing like the M-107 or M-110 is included in {\AirStr}, but again I would assert that it is tracked mainly for mobility. Effectively, these are towed artillery installed in an exposed position on a tracked hull and have almost no protection for most of the crew.

    Most tracked SAM vehicles are classified as hard targets in {\AirStr} and {\EOTG}, despite many of them having completely exposed missiles. If an air-burst bomb explodes above a SAM-4 launcher, it is not going to shrug it off. I think many SAM vehicles are tracked mainly for \emph{mobility} and perhaps for crew protection, but in terms of a mission-kill on their weapon systems, they are soft targets. I've not changed anything yet, but we need to think about this.
}

\rulesubsection{Representation}

The briefing documents show the graphical representation of ground units as counters.%
\note{Representation}{
    The representation is based on NATO symbols, including those for infantry, armor, anti-armor, air defense, guns, missiles, mortar, tube artillery, and rocket artillery. One addition is that a rectangle denotes an unarmored vehicle such as a truck or tracked carrier. A truck or mobile unit has two wheel symbols in the lower position, whereas a tracked carrier has a tracked symbol. The mortar symbol is modified to be simply an open circle, for compactness.

    All vehicular units have either an armored or unarmored vehicle symbol, all infantry units have an infantry symbol, and towed units have none of these.

    The letter T in the central position indicates a transport capacity (e.g., APCs, IFVs, cargo trucks, and tracked cargo carriers). Words and letters in the lower position indicate variants on a basic type: the letters H, M, L, and O on an armored unit indicate heavy, medium, light, and open armored protection; the word WPN or FAC on an infantry unit indicates a weapons or FAC unit; and the letter H on an infantry SAM unit indicates a heavy version of the launcher (e.g., the LML version of Starstreak).
}
The size of a unit is indicated by one, two, or three dots (for squads, sections, or platoons) or one vertical bar (for batteries) in an upper position.
Hard targets are indicated by an underlined defense strength and open targets by an underline defense strength followed by an exclamation point.

\advancedrule

\rulesubsection{Mobility}

The mobility of a ground unit restricts the hexes it can occupy according to the terrain in and around those hexes. Infantry units, tracked armored units, wheeled armored units other than armored trucks, and tracked unarmored units have unrestricted (U) mobility. Trucks, armored trucks, and other wheeled unarmored units have good (G) or poor (P) mobility. All other units have towed (T) mobility.%
\note{Mobility}{
    The inspiration for this rule was \href{https://commons.wikimedia.org/wiki/File:Aerial_view_of_the_road_through_Mitla_Pass_dotted_with_wrecked_Egyptian_vehicles_and_armour_after_Israel_air_force_attacks._June_1967._D326-013.jpg}{this photograph} of the Mitla Pass from 1967, which shows trucks destroyed on and just off a road.
    Also, in the Falklands War roads could be traversed by Landrovers and AMLs, but off-road terrain was impassable to everything except the CVR(T)s and Bv~202s despite being open.
}

A unit with unrestricted mobility may occupy: any urban hex; any clear hex; or any forest hex.
A unit with good mobility may occupy: any urban hex; any clear hex; or any forest hex with a road or trail.
A unit with poor mobility may occupy: any urban hex; any clear or forest hex with a road, trail, or runway; or any clear hex adjacent to an urban hex or another clear hex with a road, trail, or runway.
A unit with towed mobility may occupy hexes according to the mobility of its transport.

When necessary, a scenario must specify whether trucks, armored trucks, and other wheeled unarmored units have good or poor mobility, according to the terrain and the capabilities of the vehicles represented, and assign transport units to towed units.\note{Truck Mobility}{
    For example, a scenario set in the open desert of southern Iraq might specify that trucks have good mobility, but one set in the Falkland Islands might specify that they have poor mobility. Similarly, a scenario might specify that cargo trucks have poor mobility, but mobile artillery units have good mobility. Finally, a scenario might not assign an explicit transport unit to a towed artillery battery, but might state that it has been transported by helicopter and so can be placed as if it had a transport unit with good mobility.
}
A scenario may modify the mobility classes of ground units and the associated restrictions.
\note{Modifying Mobility Classes}{

    One common modification is to explicitly indicate the hexes occupied by each ground unit; in this case this rule is largely superfluous (although if we allow ground units to move, then the restrictions can still be relevant).

    Panhard AMLs are wheeled armored vehicles and so normally have unrestricted mobility. Nevertheless, in the Falklands War they were restricted to roads and trails, and this corresponds to poor mobility.
    A scenario could fix this problem by explicitly giving them poor mobility.
}

\rulesubsection{Composite Platoons}

At the start of the game, an infantry, infantry weapons, infantry mortar, or infantry HQ platoon may attach one or two infantry squads to form a composite platoon.
The attached squads are commonly infantry SAM or FAC squads.
Once formed, a composite platoon may not be split into its constituent units.%
\note{Composite Platoons}{
    This rule comes from {\TSOH}. It doesn't have counters for infantry SAMs, so in scenarios V-11 and V-22 it says treat infantry platoons as if they have an infantry SAM squad. This rule helps with stacking (since attached squads do not count for stacking) and with keeping your precious FACs and SAM teams anonymous. It also works nicely with the transport rule, since attached squads do not need their own transport but can ride with the platoon.
}

An attached squad is not represented by a separate counter, does not count for stacking, and may not be sighted, identified, or attacked directly.
Instead, it is considered to form part of the composite platoon.
An attached squad retains its normal capabilities (for example, an attached SAM squad may engage aircraft and an attached FAC squad may sight and mark targets).

% LTeX: rules-=ENGLISH_WORD_REPEAT_BEGINNING_RULE

A composite platoon has the same damage resilience, defense strength, target class, mobility, and sighting range as its constituent platoon.
If a composite platoon is sighted, no information is given on any attached squads.
If a composite platoon is identified, any attached squads are also identified (for example, “infantry platoon with an attached FAC squad”).
If a composite platoon is suppressed, each attached squad is suppressed too.
If a composite platoon suffers damage, then for each D of damage, roll one die for each attached squad and on a \minusafter{3} the attached squad is killed.
If a composite platoon is killed, each attached squad is killed too.
VPs are awarded for damage to each of the constituent units.%
\note{Damage to Attached Squads}{
    The roll of \minusafter{4} seems about right. The probability of killing an attached squad is 40\% for 1D and 64\% for 2D. Those numbers are fairly good approximations to 1/3 and 2/3. If the roll was \minusafter{3}, they would be 30\% and 51\%, and that seems a little low.

    Are the damage rolls secret?
}

% LTeX: rules+=ENGLISH_WORD_REPEAT_BEGINNING_RULE

\rulesubsection{Transport}

An infantry or towed unit may be transported by a specific associated transporting IFV, APC, armored truck, truck, or tracked-carrier unit of at least the same size.
When necessary, a scenario will establish the association between units and their specific transports.
%
\note{Transportation}{
    The idea of transporting units and cargo comes from the “Thunder and Fire” scenario from {\EOTG}. The rules here simply generalize it slightly, incorporate the idea of mobility, and add rules for mounting and dismounting.
}

A transported unit is not represented by a separate counter, does not count for stacking, and may not be sighted, identified, or attacked directly. It may not attack other ground units or aircraft, may not launch SAMs, may not use radar, and may not act as a FAC.
%Instead, a transported unit is considered to form part of the transporting unit.

% LTeX: rules-=ENGLISH_WORD_REPEAT_BEGINNING_RULE

If a transporting unit is sighted, no information is given on the transported unit. If a transporting unit is identified, no information is given on a transported infantry unit, but any transported towed unit is also identified (for example, “truck platoon towing an artillery battery” or “truck platoon towing a ZPU-4 battery”). If a transporting unit suffers damage, then its transported unit suffers the same damage.
VPs are awarded for damage to both the transporting and transported unit.

% LTeX: rules+=ENGLISH_WORD_REPEAT_BEGINNING_RULE

A unit may start the game either transported or not. If not transported, it may be placed according to either its mobility or that of its associated transport.%
\note{Mobility}{
    Being placed according to the mobility of its transport allows a tracked transport to place a towed unit in a forest hex.
}

An infantry unit in the same hex as its associated transport may declare that it is mounting the transport at the start of a Ground Unit Interaction Phase.
Both units must not be suppressed, and the transport must not have suffered more damage than the infantry unit.
At the end of the Ground Unit Interaction Phase five game turns later, the unit is considered to be transported, and its counter is removed from the map.

An infantry unit being transported may declare that it is dismounting at the start of a Ground Unit Interaction Phase.
The transport unit must not be suppressed.
At the end of the Ground Unit Interaction Phase two game turns later, the unit is considered to be no longer transported, and its counter is placed on the map in the same hex as its transport.
If placing the counter would violate stacking requirements, the unit must abort the dismounting process.

During these processes, the infantry unit may not perform any other action and the transport may not move and may not carry out ground attacks, but may perform barrage fire if it has this capability.
The infantry unit may abort either of these process at the start of any Ground Unit Interaction Phase.
If either are suppressed during the process, the unit must abort the process in the next Ground Unit Interaction Phase.%
\note{Mounting and Dismounting Infantry}{
    The rules for mounting and dismounting are completely new. I'm very open to suggestions about this.
}

A towed unit may not mount or dismount during the game.
